% Generated from locale/tr/content/incompleteness/incompleteness-provability/first-incompleteness-thm.tex; do not hand-edit.


\olfileid[tr]{inc}{inp}{1in}

\olsection{Birinci Eksiklik Teoremi}

Artık G\"odel'in birinci eksiklik teoremine ilişkin özgün ispatını
açıklayabiliriz. $\Th{T}$, aritmetik dilini genişleten bir dilde hesaplanabilir
biçimde aksiyomlaştırılmış herhangi bir kuram olsun ve $\Th{T}$,
$\Th{Q}$'nun aksiyomlarını içersin. Bu, özellikle $\Th{T}$'nin hesaplanabilir
fonksiyonları ve bağıntıları temsil ettiği anlamına gelir.

Formüllerin ve ispatların sayılarla makul bir biçimde kodlandığı
varsayıldığında $\Prf[\Th{T}](x,y)$ bağıntısının hesaplanabilir olduğunu
savunduk; burada $\Prf[\Th{T}](x,y)$, ancak ve ancak $x$, $\Th{T}$ içindeki
!!a{derivation}{}in G\"odel sayısı ve bu türetimin sonundaki !!{formula}'nın
G\"odel sayısı $y$ ise geçerlidir. Hatta G\"odel, aklındaki belirli kuram için, makalesindeki 45
fonksiyon ve bağıntı listesini kullanarak bu bağıntının ilkel özyinelemeli
olduğunu gösterebilmiştir. Kırk beşinci bağıntı $x B y$, onun seçtiği özel
$\Th{T}$ için tam olarak $\Prf[\Th{T}](x,y)$'dir. G\"odel'in makalesinde
``özyinelemeli'' sözcüğünü kullandığı yerde bugün ``ilkel özyinelemeli''
ifadesini kullanacağımızı anımsayın.

$\Prf[\Th{T}](x,y)$ hesaplanabilir olduğundan $\Th{T}$ içinde temsil
edilebilir. Onu temsil eden formülü $\OPrf[\Th{T}](x,y)$ ile göstereceğiz.
$\OProv[\Th{T}](y)$,
$\lexists[x][\OPrf[\Th{T}](x,y)]$ formülü olsun. Bu, G\"odel'in listesindeki
46'ncı bağıntı olan $\fn{Bew}(y)$'yi betimler. G\"odel'in belirttiği gibi bu,
``özyinelemeli olduğu ileri sürülemeyen'' tek bağıntıdır. Muhtemelen şunu
kastediyordu: tanımdan onun hesaplanabilir olduğu açık değildir; sonraki
gelişmeler gerçekten de hesaplanabilir olmadığını göstermiştir.

$\Th{T}$, $\Th{Q}$'yu içeren !!{axiomatizable} bir kuram olsun. O zaman
$\Prf[\Th{T}](x,y)$ karar verilebilirdir; dolayısıyla $\Th{Q}$ içinde
!!a{formula}~$\OPrf[\Th{T}](x,y)$ ile temsil edilebilir. $\OProv[\Th{T}](y)$
yukarıda betimlediğimiz formül olsun. Sabit nokta lemması uyarınca,
$\Th{Q}$'nun (dolayısıyla $\Th{T}$'nin)
\begin{equation}
\ollabel{eqn:qpf}
!G_\Th{T} \liff \lnot \OProv[\Th{T}](\gn{!G_\Th{T}}).
\end{equation}
formülünü !!{derive} ettiği bir $!G_\Th{T}$ formülü vardır. $!G_\Th{T}$'nin
özünde ``$!G_\Th{T}$, $\Th{T}$ içinde !!{derivable} değildir'' dediğine
dikkat edin.

\begin{lem}\ollabel{lem:cons-G-unprov}
$\Th{T}$, $\Th{Q}$'yu genişleten tutarlı ve !!{axiomatizable} bir kuramsa
$\Th{T} \Proves/ !G_\Th{T}$ olur.
\end{lem}

\begin{proof}
$\Th{T}$'nin $!G_\Th{T}$'yi !!{derive} ettiğini varsayın. O zaman \emph{gerçekten}
!!a{derivation} vardır ve bir $m$ sayısı için
$\Prf[\Th{T}](m,\Gn{!G_\Th{T}})$ bağıntısı geçerlidir. Fakat bu durumda
$\Th{Q}$, $\OPrf[\Th{T}](\num m,\gn{!G_\Th{T}})$ tümcesini !!{derive} eder.
Dolayısıyla $\Th{Q}$, tanım gereği $\OProv[\Th{T}](\gn{!G_\Th{T}})$ olan
$\lexists[x][\OPrf[\Th{T}](x,\gn{!G_\Th{T}})]$ formülünü !!{derive} eder.
\olref{eqn:qpf} uyarınca $\Th{Q}$, $\lnot !G_\Th{T}$'yi !!{derive} eder;
$\Th{T}$, $\Th{Q}$'yu genişlettiğinden aynısı $\Th{T}$ için de geçerlidir.
$\Th{T}$, $!G_\Th{T}$'yi !!{derive} ederse $\lnot !G_\Th{T}$'yi de
!!{derive} ettiğini ve bu nedenle tutarsız olacağını göstermiş olduk.
\end{proof}

\begin{defn}
\ollabel{thm:oconsis-q}
Bir $\Th{T}$ kuramına şu koşul sağlanıyorsa \emph{$\omega$-tutarlı} denir:
$\lexists[x][!A(x)]$ herhangi bir tümceyse ve $\Th{T}$,
$\lnot !A(\num 0)$, $\lnot !A(\num 1)$, $\lnot !A(\num 2)$, \dots
formüllerini !!{derive} ediyorsa $\Th{T}$, $\lexists[x][!A(x)]$'i ispatlamaz.
\end{defn}

Her $\omega$-tutarlı kuramın ayrıca tutarlı olduğuna dikkat edin. Bunun nedeni,
$\Th{T}$ tutarsızsa her $!A$ için $\Th{T} \Proves !A$ olmasıdır. Özellikle
$\Th{T}$ tutarsızsa her $n$ için hem $\lnot !A(\num n)$'yi !!{derive} eder
hem de $\lexists[x][!A(x)]$'i !!{derive} eder. Dolayısıyla $\Th{T}$ tutarsızsa
$\omega$-tutarsızdır. Karşıt tersinden, $\Th{T}$ $\omega$-tutarlıysa tutarlı
olmak zorundadır.

\begin{lem}\ollabel{lem:omega-cons-G-unref}
$\Th{T}$, $\Th{Q}$'yu genişleten $\omega$-tutarlı ve !!{axiomatizable} bir
kuramsa $\Th{T} \Proves/ \lnot !G_\Th{T}$ olur.
\end{lem}

\begin{proof}
$\Th{T}$, $\lnot !G_\Th{T}$'yi !!{derive} ederse $\omega$-tutarsız olduğunu
göstereceğiz. $\Th{T}$'nin $\lnot !G_\Th{T}$'yi !!{derive} ettiğini varsayın.
$\Th{T}$ tutarsızsa $\omega$-tutarsızdır ve işimiz biter. Aksi durumda
$\Th{T}$ tutarlıdır; bu yüzden \olref{lem:cons-G-unprov} uyarınca
$!G_\Th{T}$'yi !!{derive} etmez. $\Th{T}$ içinde $!G_\Th{T}$'nin
!!{derivation}{}i bulunmadığından $\Th{Q}$
\[
\lnot \OPrf[\Th{T}](\num 0, \gn{!G_\Th{T}}), \lnot \OPrf[\Th{T}](\num 1,
\gn{!G_\Th{T}}), \lnot \OPrf[\Th{T}](\num 2, \gn{!G_\Th{T}}), \dots
\]
formüllerini !!{derive} eder; dolayısıyla $\Th{T}$ de eder. Öte yandan
\olref{eqn:qpf} uyarınca $\lnot !G_\Th{T}$,
$\lexists[x][\OPrf[\Th{T}](x,\gn{!G_\Th{T}})]$ formülüne denktir. Öyleyse
$\Th{T}$ $\omega$-tutarsızdır.
\end{proof}

\begin{prob}
  Her $\omega$-tutarlı kuram tutarlıdır. Tersinin geçerli olmadığını, yani
  tutarlı fakat $\omega$-tutarsız kuramlar bulunduğunu gösterin. Bunu
  $\Th{Q} \cup \{\lnot !G_\Th{Q}\}$ kuramının tutarlı fakat
  $\omega$-tutarsız olduğunu göstererek yapın.
\end{prob}

\begin{thm}
\ollabel{thm:first-incompleteness} $\Th{T}$, $\Th{Q}$'yu genişleten herhangi
bir $\omega$-tutarlı, !!{axiomatizable} kuram olsun. O zaman $\Th{T}$ tam
değildir.
\end{thm}

\begin{proof}
  $\Th{T}$ $\omega$-tutarlıysa tutarlıdır; dolayısıyla
  \olref{lem:cons-G-unprov} uyarınca $\Th{T} \Proves/ !G_\Th{T}$ olur.
  \olref{lem:omega-cons-G-unref} uyarınca da
  $\Th{T} \Proves/ \lnot !G_\Th{T}$ olur. Bu, $\Th{T}$ ne $!G_\Th{T}$'yi
  ne de $\lnot !G_\Th{T}$'yi !!{derive} ettiğinden eksik olduğu anlamına
  gelir.
\end{proof}

